\section{Machine-Checked Companion}\label{app:formalization}

\noindent
Nothing in this appendix is used by the mathematics above.  Every definition,
theorem and proof in the body is complete in ordinary mathematical language and
should be read and judged as such.  The development recorded here is a
correctness check on the author's own reasoning, not evidence offered to the
reader, and it establishes neither significance nor priority.  It is reported
because it exists, and because its trust boundary should be stated explicitly
rather than left implicit.

The repository pins Lean~4.33.0.  The proof package contains a small
publication-facing facade in \texttt{ResolutionMasterTheorems.lean} that
exposes the structural results used in the reorganized paper.  In particular,
the following declarations are checked directly:
\begin{lstlisting}[style=lean]
#check ResolutionSemantics.MasterTheorems.relativeFinitePatternRealization
#check ResolutionSemantics.MasterTheorems.orbitCompressionCauchy
#check ResolutionSemantics.MasterTheorems.finitePatternEscapeNoGeneratedLimit
#check ResolutionSemantics.MasterTheorems.compressedEscapingOrbitNotComplete
#check ResolutionSemantics.MasterTheorems.compressedEscapingOrbitEmbeddingNotSurjective
#check ResolutionSemantics.MasterTheorems.finiteBasePropernessCriterion
#check ResolutionSemantics.MasterTheorems.oldFixingPropernessViaCombinedMaster
#check ResolutionSemantics.MasterTheorems.oldFixingRanksUnbounded
\end{lstlisting}

The finite-pattern result is implemented in
\texttt{ResolutionFinitePatternRealization.lean}.  The trajectory-compression
abstraction is in \texttt{ResolutionOrbitCompressionMaster.lean}; its finite
code applies only to the relevant deterministic orbit and does not require a
finite ambient observer carrier.  The anti-limit theorem is implemented in
\texttt{ResolutionFinitePatternAntiLimit.lean}; it formalizes the selected
candidate/prefix construction, overflow entry, and overflow absorption.  Both
the finite-base and old-fixing properness theorems are re-derived as instances
of the combined theorem.

The command
\begin{lstlisting}[style=lean]
bash scripts/verify.sh
\end{lstlisting}
checks the manuscript-linked public API, compiles the proof dependency chain,
rejects direct proof holes or local axioms in the paper sources, and runs the
axiom audit.  The accepted theorem dependencies are restricted to
\texttt{propext}, \texttt{Classical.choice}, and \texttt{Quot.sound}.

The formalization uses classical finite indexing and witness selection in
explicit places.  The resulting trust claim is therefore relative to Lean's
kernel, the imported standard library, and the ordinary classical principles
reported by the audit.  No foundation-independent consistency claim is made.

\subsection{Theorem map}

The declarations corresponding to the numbered results of the body are
listed below, grouped by section.  They are provided for traceability only.

\paragraph{Free completion substrate (\S
ef{sec:free-completion})}
\begin{itemize}\setlength{\itemsep}{0pt}
  \item \leanname{ResolutionSemantics.freeCompletionUniversal}
  \item \leanname{ResolutionSemantics.expressionKernelQuotientInjective}
  \item \leanname{ResolutionSemantics.expressionKernelQuotientSurjective}
\end{itemize}

\paragraph{Observers and finite-pattern realization (\S
ef{sec:finite-separation})}
\begin{itemize}\setlength{\itemsep}{0pt}
  \item \leanname{ResolutionSemantics.MasterTheorems.relativeFinitePatternRealization}
  \item \leanname{ResolutionSemantics.finiteSeparationBound}
  \item \leanname{ResolutionSemantics.MasterTheorems.relativeFiniteComplementSeparation}
  \item \leanname{ResolutionSemantics.RelativeObservers.finiteRelativeComplement_not_closed_under_binary_products}
\end{itemize}

\paragraph{Completion, compression, escape (\S
ef{sec:completion})}
\begin{itemize}\setlength{\itemsep}{0pt}
  \item \leanname{ResolutionSemantics.MasterTheorems.orbitCompressionCauchy}
  \item \leanname{ResolutionSemantics.MasterTheorems.finitePatternEscapeNoGeneratedLimit}
  \item \leanname{ResolutionSemantics.MasterTheorems.compressedEscapingOrbitEmbeddingNotSurjective}
  \item \leanname{ResolutionSemantics.MasterTheorems.compressedEscapingOrbitNotComplete}
  \item \leanname{ResolutionSemantics.MasterTheorems.finiteBasePropernessCriterion}
  \item \leanname{ResolutionSemantics.MasterTheorems.oldFixingPropernessViaCombinedMaster}
  \item \leanname{ResolutionSemantics.MasterTheorems.oldFixingRanksUnbounded}
  \item \leanname{ResolutionSemantics.completionComplete}
  \item \leanname{Resolution.External.FilteredTotalAlg.completedResolutionFilteredAlgebra\_initial}
  \item \leanname{Resolution.Probe.depthCauchy\_not\_observationalCauchy}
  \item \leanname{Resolution.Probe.comb\_separatesAt\_one}
  \item \leanname{Resolution.Probe.NatProbe.natComb\_pairwise\_separated\_at\_zero}
\end{itemize}

\paragraph{Equational conservativity (\S
ef{sec:equations})}
\begin{itemize}\setlength{\itemsep}{0pt}
  \item \leanname{ResolutionSemantics.equationConservative}
\end{itemize}

\paragraph{Arithmetic instances (\S
ef{sec:arithmetic})}
\begin{itemize}\setlength{\itemsep}{0pt}
  \item \leanname{ResolutionSemantics.NatDivision.singularFamilyInjective}
  \item \leanname{ResolutionSemantics.NatDivision.singularFamilyDisjoint}
  \item \leanname{ResolutionSemantics.NatDivision.oldFixingCriterion}
  \item \leanname{ResolutionSemantics.NatDivision.separationRanksUnbounded}
  \item \leanname{ResolutionSemantics.IntDivision.singularFamilyInjective}
  \item \leanname{ResolutionSemantics.IntDivision.zeroDivZeroNotOld}
  \item \leanname{ResolutionSemantics.IntDivision.completionEmbeddingNotSurjective}
  \item \leanname{ResolutionSemantics.IntDivision.separationRanksUnbounded}
\end{itemize}

\paragraph{Many-sorted localization (\S
ef{sec:diaconescu-bridge})}
\begin{itemize}\setlength{\itemsep}{0pt}
  \item \leanname{Resolution.Diaconescu.ManySorted.term\_evalPartial\_sound}
  \item \leanname{Resolution.Diaconescu.ManySorted.term\_evalPartial\_complete}
  \item \leanname{Resolution.Diaconescu.ManySorted.alpha\_satisfaction\_condition}
  \item \leanname{Resolution.Diaconescu.ManySorted.gamma\_satisfiesGamma}
  \item \leanname{Resolution.Diaconescu.ManySorted.betaGammaPartialAlgEquiv}
  \item \leanname{Resolution.Diaconescu.ManySorted.gammaBetaHomEquiv}
  \item \leanname{Resolution.Diaconescu.ManySorted.gamma\_preserves\_initiality}
  \item \leanname{Resolution.Diaconescu.ManySorted.beta\_of\_initial\_encoded\_is\_initial}
  \item \leanname{Resolution.Diaconescu.ManySorted.semantic\_consequence\_equivalence}
  \item \leanname{Resolution.Diaconescu.InitialExistence.quasiTheory\_has\_initial\_partial\_model}
  \item \leanname{Resolution.Diaconescu.BinarySpecialization.generatedEquiv}
  \item \leanname{Resolution.Diaconescu.BinarySpecialization.gammaEquiv}
\end{itemize}
